% Audited unnumbered auxiliary.  The fixed proper-family obstruction remains.
\paragraph{The rational exceptional spectrum of the ordered tetrahedron
(audited auxiliary calculation).}
Let $k_0=\overline{\mathbf F}_p$, and let $E_0/k_0$ be the reduced
exceptional surface in the ordered resolution TG.1--TG.5.  Thus its
four components are the blowups of $\mathbf P^2$ at $15,10,5,0$
nonvertex edge points, respectively.  In this paragraph the subscript
$0$ denotes the field of definition, not the unresolved first blowup
used in TG.2.  For every quasi-compact separated $k_0$-scheme $P$ for
which the indicated products are Noetherian, the external-product map
is a natural equivalence of rational $K(P)$-module spectra
\begin{equation}
 K(E_0)_{\mathbf Q}\otimes_{K(k_0)_{\mathbf Q}}K(P)_{\mathbf Q}
 \ \xrightarrow{\ \sim\ }\ K(E_0\times_{k_0}P)_{\mathbf Q}.
 \tag{TS.1}
\end{equation}
Here $K$ is ordinary Bass-completed $K$-theory of perfect complexes.
There is a noncanonical, parameter-independent splitting
\begin{equation}
 K(E_0\times_{k_0}P)_{\mathbf Q}
 \simeq K(P)_{\mathbf Q}^{\oplus33}
       \oplus \Sigma^{-2}K(P)_{\mathbf Q}.
 \tag{TS.2}
\end{equation}
The equivalence (TS.1), unlike the displayed choice of splitting, is
natural in $P$.  No equivalence with a base change of $K(Z)$ is asserted.

\emph{Step 1: a comparison of the entire restriction diagrams.}
For each nonempty intersection $E_I$ of components, consider its
canonical external-product map
\[
 K(E_I)_{\mathbf Q}\otimes_{K(k_0)_{\mathbf Q}}K(P)_{\mathbf Q}
       \longrightarrow K(E_I\times P)_{\mathbf Q}.
\]
It is an equivalence.  For projective spaces and points this follows
from the projective-bundle formula.  For the blown-up planes it follows
successively from the blowup formula in a smooth point.  The formulas
are induced by actual pullback, tensor, and perfect proper pushforward
functors; hence their comparison with the external product is the
projection formula, not a comparison of ranks.  We use
\citet[Corollaries 3.6 and 4.4, author pp.~12 and 16]
{D_Khan2020Blowup} only for these ordinary regular projective bundles
and blowups.

The comparison maps themselves commute with all restriction maps and
all their coherences, because external product is natural before
choosing any projective-bundle bases.  Thus there is an objectwise
equivalence of the full finite cubical restriction diagrams.  This
avoids choosing incompatible cellular flags or discarding gluing
automorphisms at the nodes.

\emph{Step 2: descent for the closed union.}
For a binary closed union in characteristic $p$, the local square is a
Milnor square of rings.  Theorem C of
\citet[p.~30]{rootGeisserHesselholt2011} makes its birelative spectrum
zero after inverting $p$, and finite affine Mayer--Vietoris globalizes
this assertion.  Nilpotent reductions are harmless by their Theorem A.
Iterating over the four components and their intersections therefore
expresses $K(E_0)_{\mathbf Q}$ and $K(E_0\times P)_{\mathbf Q}$ as
the finite homotopy limits of their closed-cover diagrams.  Tensor
product of module spectra is exact and hence commutes with finite
homotopy limits.  Applying it to the objectwise equivalence in Step~1
proves (TS.1).  This argument uses ordinary $K$ after rationalization;
it does not substitute $KH$ without a comparison.

\emph{Step 3: the finite rational complex over $k_0$.}
The rational $K$-spectrum of $k_0$ is $H\mathbf Q$: the positive
groups are torsion and the negative groups of a field vanish.
For the positive statement one may use Quillen's finite-field
calculation and filtered colimits; the exact displayed calculation
is recorded in \citet[Chapter~IV, Corollary~1.13 and Example~2.6]
{D_Weibel2013KbookIV}.  Only rational vanishing is required here.
Consequently every stratum's rational spectrum has homotopy solely
in degree zero.  Its canonical identification with the Eilenberg--
MacLane spectrum of its $K_0$ is natural in all restriction maps.
The entire diagram, including its higher coherences, can therefore be
computed in rational vector spaces.  There are no discarded positive
$K_1(k_0)_{\mathbf Q}$ entries.

The normalized cochain complex is
\begin{equation}
 \bigoplus_{i=0}^3 K_0(E_i)_{\mathbf Q}
   \longrightarrow \bigoplus_{i<j}K_0(E_i\cap E_j)_{\mathbf Q}
   \longrightarrow \bigoplus_{i<j<l}K_0(E_i\cap E_j\cap E_l)_{\mathbf Q}.
 \tag{TS.3}
\end{equation}
Rational Chern character on these smooth rational strata splits this
complex into its codimension-zero, one, and two parts, compatibly
with pullback.  The codimension-zero part is the ordinary cochain
complex of the boundary of the tetrahedron.  Its cohomology is
$\mathbf Q$ in degrees zero and two and is zero in degree one.

The codimension-one part is the restriction-degree map
\[
 \bigoplus_i\operatorname{Pic}(E_i)_{\mathbf Q}
       \longrightarrow \mathbf Q^{6}.
\]
It is surjective: on each edge choose one of the five point-exceptional
curves on the earlier endpoint component.  Its intersection with that
edge's strict transform is one and its intersection with every other
edge is zero, because its centre is a nonvertex point.  Its Picard
class therefore supplies the corresponding signed standard basis
vector.  The source has rank $4+30=34$, so the kernel has rank~28.
The codimension-two part is $\mathbf Q^4$ in degree zero, since each
component has $CH_0=\mathbf Z$ and the edges have no codimension-two
Chow group.  We conclude
\[
 \pi_0K(E_0)_{\mathbf Q}=\mathbf Q^{33},\qquad
 \pi_{-1}K(E_0)_{\mathbf Q}=0,\qquad
 \pi_{-2}K(E_0)_{\mathbf Q}=\mathbf Q,
\]
with all other rational groups zero.  A finite complex of rational
vector spaces splits as its cohomology.  Choosing that splitting
once over $k_0$, then applying (TS.1), proves (TS.2).

\emph{Step 4: the exact residual coefficient, and what is not proved.}
For $P=\operatorname{Spec}k$, where $k/k_0$ is algebraically closed,
the summand in degree zero arising from the final term of (TS.2) is
an actual ordinary $K_2(k)_{\mathbf Q}$ coefficient.  A generator
$\kappa\in K_{-2}(E_0)_{\mathbf Q}$ has Adams weight zero, and this
coefficient is its product with $K_2(k)$, which has Adams weight two.
The remaining weight-two part is the four component-degree classes.
For a smooth proper curve $P/k$, the same chosen summand is
$K_2(P)_{\mathbf Q}$ and evaluation at a point is the ordinary
$K_2$ specialization.  This parameter compatibility follows from
(TS.1), not just from the dimensions of (TS.3).

The product restriction $K(Z\times P)\to K(E\times P)$ has not been
computed by this argument.  Its precise characteristic-$p$ cokernel
is FF.1, and on the displayed coefficient its boundary is FF.3.
There is an important separate point-parameter conclusion:
\citet[Theorem~3.7 and Remark~(ii), author p.~8]{frDrinfeld2004Tate}
give $K_{-1}(R)=0$ for every commutative henselian local ring.
Thus, for the strict-local version of TG, FF.1 with a point parameter
does prove that every individual coefficient lifts after inverting
$p$.  The KST sequence also proves injectivity: $K_0(R)=\mathbf Z$
maps by the identity to $K_0(k)=\mathbf Z$, so a pair $(z,0)$ in its
diagonal image must be zero.  Hence $K_0(Z)[1/p]\simeq K_0(E)[1/p]$.
In particular the point coefficient product with $\partial\kappa$
vanishes in degree $-1$ after strict henselization.  This corrects
the earlier prospective point-faithfulness test.

The fixed proper product $\operatorname{Spec}R\times P$ is not a
henselian local scheme.  Drinfeld's local theorem does not make its
$K_{-1}$ vanish.  Neither (TS.2) nor the point isomorphism just proved
lifts the prescribed whole proper-curve $K_2$ class to $Z\times P$.
That full-family boundary remains outside this calculation.
The subsequent endpoint theorem EL computes the ordinary quotient
for this tetrahedral example without asserting that the prescribed
whole perfect family lifts.
